\section{Tests de validation}
Il est important de pouvoir tester son code et nous l'expliquons succintement dans cette partie.

\subsection{Tests structurels}

Plusieurs test ont été réalisés sur les clients pour vérifier le bon respect des règles du jeu.%huyjtfrgklggujh(Camille)
D'autres tests succints ont été réalisés sur certaines fonctions.

La meilleure façon de valider le fonctionnement de nos clients reste de les faire s'affronter et d'observer le graphe de jeu obtenu dont nous discutons dans la prochaine partie.

\subsection{Affichage des parties}

Un autre moyen de tester le bon fonctionnement de nos clients ou du serveur est simplement d'observer le graphe du jeu.
Dans ce sens il nous a été nécessaire de créer une fonctions d'affichage des graphes avec les joueurs.

C'est ainsi que l'on a utilisé \textit{DOT}, un language de description de graphe.
Ce language nous permet de générer un graphe pour chaque état de la partie, le tout dans un pdf.
C'est le but de la fonction \textit{graph\_dot} qui a été apercu précédemment mais non expliqué.
Celle-ci génére un fichier \textit{DOT} qui est un simple fichier texte décrivant le graphe, ce fichier est ensuite compilé en pdf en utilisant \textit{neato}.
Il nous a également fallu créer un script bash pour lier chaque pdf de chaque graphe.
Ainsi, on obtient un pdf du nombre de page correspondant au nombre de coup de la partie.

\begin{figure}[H]
  \centering
  \includegraphics[width=0.8\textwidth]{dot.png}
  \caption{Affichage du plateau de jeu}
  \label{fig:dot}
\end{figure}


La figure \ref{fig:dot} montre un début de partie avec 4 objectifs, le joueur noir sur le sommet 0 et le joueur blanc sur le sommet 90. 

\section{Conclusion}

Ce projet fut vaste, à la fois sur les plans du développement logiciel, de l’algorithmie et de la théorie. Nous avons appris à gérer des serveurs, à créer des bibliothèques dynamiques.
Nous avons aussi appris à prendre du recul sur un problème afin de créer des outils nous permettant de résoudre ces mêmes problèmes.
Enfin, nous avons dû concevoir une stratégie, la tester, la comparer à d'autres, la remettre en question, puis la peaufiner.
En résumé, le projet Coursidor fut enrichissant à bien des égards pour les jeunes programmeurs que nous sommes.

